\subsection{PACE Process for the Allocation of Preferential Seats}\label{sec:PACEadmissionprocess}
The PACE and regular applications must be submitted to the centralized system during the same time window. For the cohort included in this study, the PACE and regular admission processes were separate. Students could submit a PACE application list and, separately, a regular application list. In each list a student could include up to ten programs (i.e., college and major combinations), potentially different between the two lists. The two processes were entirely independent, and a student could obtain two admissions simultaneously, one from each process. Here we describe the PACE application and admission process.\par 
For each program listed in their PACE applications, applicants receive a distinct application score, called {\itshape Puntaje de Postulaci\'{o}n PACE} (PPP). The score is calculated based on the applicant's GPA during the four years of high school and attendance during the $11^{th}$ and $12^{th}$ grades. To reduce the occurrence of identical scores across applicants, the score is adjusted for each program, taking into account the program's geographic location and its positional ranking within the applicant's list of preferences.\footnote{The formula is $PPP=(0.8 * PRN + 0.2 * GPA)\cdot (1+bonus_{attendance}+bonus_{geog})+bonus_{listrank}$, where $PRN$ is the {\itshape Puntaje Ranking de Notas} (PRN) used to identify the top $15\%$ of students, which is based on the high school GPA with some adjustments (see Appendix \ref{sec:ptjeranking}), and $GPA$ is the raw high school GPA. The correlation coefficient between the raw GPA and the PRN is around 98\%.  The bonus for attendance rewards high school attendance and it reaches a maximum of 5\% for students who did not miss a single day and drops to 0 for those missing 15\% or more days. A bonus for geographic location of 3.5\% (5\%) is awarded for applications to universities in the same area of Chile (North, Center, or South) (region of Chile) as the student's high school, and the bonus for the rank of the program within the applicant list decreases with the program's rank; it is measured in score points, and it is 25 for applications listed first, typically representing less than 5\% of the total score, and 0 for applications listed tenth. \label{foot:PPPformula} }\par 

Applicants to each program are ranked in descending order based on their application score, and available PACE slots are allocated according to this sequence. Should the number of applicants exceed the available slots, those not immediately admitted are placed on a waiting list for their first-choice program. Subsequently, these candidates are considered for admission to the programs listed as their second choice, following the same order-based allocation process. This procedure is iteratively applied to applicants' subsequent choices. Once an applicant is accepted into a program, they are automatically withdrawn from consideration for any programs ranked lower on their preference list. This measure ensures that no applicant is admitted to more than one program. However, applicants remain eligible for programs ranked higher than their successful application, should they be initially placed on a waiting list for such programs. In instances where a student eligible for a guaranteed PACE slot fails to secure admission in any of their listed preferences, the Ministry of Education employs a proprietary algorithm to determine their placement.\par 